\section{The Heart Laid Open}

The Lord will bring to light what is hidden and disclose the counsels of hearts. St.~Paul says this while speaking of the stewardship entrusted to the apostles. He does not regard human approval as the final court, and he adds that an untroubled self-assessment does not justify him: the Lord judges (1~Cor 4:1--5). Conscience matters profoundly, but our present knowledge of ourselves is incomplete. We can miss motives, excuse habits, forget duties, or attach excessive guilt to something that was never freely chosen. Christ's judgment reaches the truth beyond these distortions.

Augustine reads the opened books of Revelation~20 as the disclosure of God's commands and of the deeds done or omitted in response. He understands the book of each life through the divine power which brings its acts before conscience (\emph{City of God} XX.14). On his account, disclosure depends upon God's knowledge of the life. A forgotten act remains known to him; an excuse we have repeated for years remains subject to his judgment. This is a reason to examine what we freely chose while we can still repent of it.

Such truth concerns interior acts as well as exterior behavior. An outwardly generous gift can become an instrument of domination; an apparently gracious silence can conceal consent to slander; diligent service can be mingled with a hunger for recognition. These mixtures call for purification. They do not require us to deny every good in an imperfect work, but they ask us to surrender the part we keep defending. The Lord can distinguish the love that was real from the vanity that attached itself to it. Examination becomes more honest when we ask him to do that work within us now.

There are also hidden victories. Restraining a cruel reply may require a struggle no one else sees. Fulfilling a duty through grief may take greater effort than the same task required before the loss. Cyril tells his hearers that their prayer, alms, and fidelity are recorded before God (XV.23). We can ask for help when a duty becomes difficult and continue it without making another person's recognition a condition of our service.

God's knowledge includes the unequal histories in which freedom is exercised. Ignorance, fear, compulsion, habit, and other psychological or social factors can diminish responsibility, sometimes even nullifying it (CCC 1735). The objective wrongness of an act and the personal culpability of its agent are distinct. The Judge knows every relevant circumstance without confusion. He knows the wound we did not choose, the light we truly lacked, the assistance we received, and the occasions on which we freely resisted what we knew to be good. No human estimate can achieve this completeness.

That truth should make our examination humble in two directions. We should not excuse a known sin by assembling a life story in which every wrong becomes inevitable. Nor should we claim a certainty of guilt which ignores involuntary suffering and the limits of freedom. A history can explain why a temptation is strong without making consent to it good. A strong emotion can accompany a holy act without becoming a sin. We bring the actual act, intention, circumstances, and degree of consent to a conscience formed by truth, receiving pastoral help where our judgment has become confused.

Conscience itself needs formation. It must be obeyed in its certain judgment, yet it can be mistaken, and negligence in seeking the truth can be culpable. Invincible ignorance has a different moral character from a refusal to learn what one should know (CCC 1790--1794). A person may carefully avoid a teaching because learning it would disturb a profitable arrangement. Another may sincerely struggle with a question and seek correction. Both can say that they followed conscience, but the phrase does not make the two histories identical. God sees the willingness or unwillingness to receive light.

Paul's warning in Romans makes the danger sharper. A person condemns another while doing the same things, as though the ability to recognize evil gave immunity from its judgment (Rom 2:1--5). The apostle addresses the judge who has turned a true moral judgment into protection for an unrepentant self. We can fall into that position through indignation which never returns to our own conduct. Another person's failure becomes absorbing; our own begins to seem small beside it. The truth we see outside ourselves should instead make us more willing to receive its searching light within.

Paul then asks what we have made of God's patience. Its purpose is to lead us to repentance. We can misread the absence of immediate punishment as approval, or take repeated opportunities to return as evidence that returning can safely be postponed. The years in which a hidden fault has gone undiscovered are still years lived before God. His kindness has accompanied them with a call to conversion. When conscience awakens after a long evasion, we may receive that awakening as mercy, even though acknowledging the truth is painful.

The same chapter places conscience within God's impartial judgment. Those who possess the revealed law are accountable for what they have received; those without that privilege still bear witness to a law written in their hearts. Their thoughts accuse or defend them, and God judges human secrets through Jesus Christ (Rom 2:12--16). We may know another person chiefly by a religious or social description. God knows that person's interior response to the truth received. He judges our response with the same completeness.

For us, this demands a distinction between hearing the truth and doing it. A Catholic can have heard a commandment many times while continuing to arrange life around its evasion. The familiarity of the words may even make their demand harder to hear. Paul returns us to obedience: the law is given to govern conduct. We ought to be grateful when a Gospel already known suddenly becomes uncomfortable, because its discomfort may reveal the particular surrender we have resisted. The judgment of conscience can then become a threshold of repentance, through which we return to Christ's judgment in the sacrament of mercy.

Luke places the demand for an account beside the gifts entrusted to the servant: more is required where more has been given (Lk 12:47--48). We should apply this first to ourselves. The faith we were taught, the sacraments available to us, the patient priest, the faithful parent, the friend who corrected us, the leisure for prayer, and the opportunities for service are gifts. They call for gratitude and a response. We cannot use them merely to compare ourselves favorably with those who received less help. Religious knowledge increases the obligation to let known truth govern life.

Cyril's warning about the secret heart ends by directing his hearers toward faithful action. No absence of witnesses hides a deed from God; the conscience will bear witness; therefore baptismal cleansing calls for a life kept clean and a fruitful use of what God entrusts (XV.25--26). We can act on this counsel even while some motives remain confused. A known sin can be confessed and a recognized excuse relinquished. Further examination, with sound pastoral counsel, can address the uncertainty which remains.

